\documentclass[a4paper,12pt]{article}

% ------------------------------------------------
% Кодировка, язык, математика
% ------------------------------------------------
\usepackage[T2A]{fontenc}
\usepackage[utf8]{inputenc}
\usepackage[russian]{babel}
\usepackage{amsmath,amssymb,amsthm,mathtools}
\usepackage{icomma}

% ------------------------------------------------
% Типографика
% ------------------------------------------------
\usepackage{helvet}
\renewcommand{\familydefault}{\sfdefault}
\usepackage{microtype}
\usepackage{setspace}
\onehalfspacing
\usepackage{parskip}

% ------------------------------------------------
% Страница
% ------------------------------------------------
\usepackage[
  a4paper,
  left=20mm,
  right=20mm,
  top=20mm,
  bottom=22mm
]{geometry}

% ------------------------------------------------
% Графика
% ------------------------------------------------
\usepackage{tikz}
\usetikzlibrary{calc}
\usetikzlibrary{arrows.meta,positioning,shapes.geometric}
\usepackage{tkz-euclide}

\usepackage{pgfplots}
\pgfplotsset{compat=1.18}

% ------------------------------------------------
% Таблицы и списки
% ------------------------------------------------
\usepackage{tabularx,booktabs,longtable}
\usepackage{enumitem}

% ------------------------------------------------
% Служебные пакеты
% ------------------------------------------------
\usepackage{xcolor}
\usepackage{hyperref}
\usepackage{fancyhdr}

% ------------------------------------------------
% Палитра
% ------------------------------------------------
\definecolor{accent}{HTML}{008080}
\definecolor{darkaccent}{HTML}{004D4D}
\definecolor{textgray}{HTML}{333333}
\definecolor{softgray}{HTML}{777777}
\definecolor{lightaccent}{HTML}{DCEEEE}

\color{textgray}

% ------------------------------------------------
% Настройки ссылок
% ------------------------------------------------
\hypersetup{
  colorlinks=true,
  linkcolor=darkaccent,
  urlcolor=darkaccent,
  pdfborder={0 0 0}
}

% ------------------------------------------------
% Колонтитулы
% ------------------------------------------------
\pagestyle{fancy}
\fancyhf{}
\fancyhead[L]{\small\color{softgray}Стереометрия}
\fancyhead[R]{\small\color{softgray}Ксения Клыкова}
\fancyfoot[C]{\small\color{softgray}\thepage}
\renewcommand{\headrulewidth}{0.3pt}
\renewcommand{\headrule}{%
  \hbox to\headwidth{\color{lightaccent}\leaders\hrule height \headrulewidth\hfill}}
\setlength{\headheight}{15pt}

% ------------------------------------------------
% TkZ
% ------------------------------------------------
\tkzSetUpLine[line width=1pt,color=accent]
\tkzSetUpPoint[color=accent,fill=accent,size=3]
\tkzSetUpLabel[color=black,font=\small]

% ------------------------------------------------
% Блок-схемы
% ------------------------------------------------
\tikzset{
  flowstart/.style={
    ellipse,
    draw=accent,
    line width=0.9pt,
    align=center,
    minimum width=42mm,
    minimum height=10mm,
    inner sep=4pt
  },
  flowaction/.style={
    rectangle,
    rounded corners=3pt,
    draw=accent,
    line width=0.9pt,
    align=center,
    text width=93mm,
    minimum height=12mm,
    inner sep=5pt
  },
  flowdecision/.style={
    diamond,
    aspect=3.3,
    draw=darkaccent,
    line width=0.9pt,
    align=center,
    text width=52mm,
    inner sep=2pt
  },
  flowarrow/.style={
    -{Latex[length=3mm,width=2mm]},
    line width=0.9pt,
    color=darkaccent
  }
}

% ------------------------------------------------
% Пользовательские команды
% ------------------------------------------------
\newcommand{\lessonsection}[1]{%
  \vspace{0.4em}
  {\Large\bfseries\color{darkaccent}#1}\par
  \vspace{0.25em}
  {\color{lightaccent}\rule{\linewidth}{0.8pt}}
  \vspace{0.6em}
}

\newcommand{\key}[1]{\textbf{\color{darkaccent}#1}}

\begin{document}

% ================================================================
% ТИТУЛЬНЫЙ ЛИСТ
% ================================================================
\begin{titlepage}
\thispagestyle{empty}

\begin{center}

\vspace*{24mm}

{\large\color{softgray}Математика $\cdot$ 10 класс}

\vspace{12mm}

{\Huge\bfseries\color{darkaccent}Чек-лист}

\vspace{5mm}

{\LARGE\bfseries
Параллельность в пространстве\\[2mm]
и переход к плоской задаче}

\vspace{13mm}

{\large
Прямые и плоскости $\cdot$
средняя линия $\cdot$
подобие треугольников}

\vfill

{\large \today}

\vspace{9mm}

{\large
Лёвин Артём Александрович\\[2mm]
\textbf{эксклюзивно для Ксении Клыковой}}

\vspace{20mm}

\begin{tikzpicture}[remember picture]
\draw[
  accent,
  line width=1.2pt
]
(-6.6,0)
.. controls (-4.2,1.0) and (-2.4,-0.85) ..
(0,0)
.. controls (2.4,0.85) and (4.3,-1.0) ..
(6.6,0);
\end{tikzpicture}

\vspace{5mm}

\end{center}
\end{titlepage}

\setcounter{page}{1}

% ================================================================
% ОСНОВНОЙ ЧЕК-ЛИСТ
% ================================================================
\lessonsection{1. Что важно видеть в пространственной задаче}

Перед доказательством полезно разделить информацию на четыре типа:

\begin{itemize}[leftmargin=8mm,itemsep=2mm]
  \item какие \key{точки принадлежат прямым};
  \item какие \key{точки и прямые принадлежат плоскости};
  \item какие прямые \key{параллельны};
  \item какие объекты \key{пересекаются}.
\end{itemize}

В стереометрии рисунок помогает увидеть конструкцию, однако утверждать можно
только то, что следует из условия или уже доказано.

\medskip

\key{Ключевой принцип занятия:}

\[
\boxed{\text{пространственную задачу часто полезно свести к треугольнику}}
\]

После выбора подходящей плоскости становятся доступны знакомые планиметрические
инструменты: средняя линия, подобие, отношения отрезков.

% ----------------------------------------------------------------
\lessonsection{2. Две точки прямой в плоскости}

\key{Свойство.}
Если две различные точки прямой принадлежат плоскости, то вся прямая
принадлежит этой плоскости.

Пусть
\[
A\in c,\qquad B\in c,\qquad A\in\alpha,\qquad B\in\alpha,
\qquad A\ne B.
\]
Тогда
\[
c\subset\alpha.
\]

Это один из самых коротких и полезных способов доказать принадлежность прямой
плоскости.

\subsection*{\color{darkaccent}Типовая ситуация}

Пусть
\[
a\parallel b,\qquad a\subset\alpha,\qquad b\subset\alpha,
\]
а прямая \(c\) пересекает \(a\) и \(b\):
\[
A=c\cap a,\qquad B=c\cap b.
\]

Тогда
\[
A\in\alpha,\qquad B\in\alpha.
\]

Поскольку \(a\parallel b\), точки \(A\) и \(B\) различны. Следовательно,

\[
\boxed{c\subset\alpha}.
\]

\begin{center}
\begin{tikzpicture}[scale=0.9]
  \tkzDefPoint(0,0){P}
  \tkzDefPoint(7,0.7){Q}
  \tkzDefPoint(8.3,4.2){R}
  \tkzDefPoint(1.3,3.5){S}

  \draw[fill=lightaccent,line width=0.7pt,draw=accent]
    (P)--(Q)--(R)--(S)--cycle;

  \tkzDefPoint(1.4,1.15){A1}
  \tkzDefPoint(6.8,1.7){A2}
  \tkzDefPoint(2.1,2.75){B1}
  \tkzDefPoint(7.5,3.3){B2}

  \tkzDrawSegment(A1,A2)
  \tkzDrawSegment(B1,B2)

  \tkzDefPoint(3.15,1.33){A}
  \tkzDefPoint(5.45,3.09){B}
  \tkzDrawSegment(A,B)

  \tkzDrawPoints(A,B)
  \tkzLabelPoints[below](A)
  \tkzLabelPoints[above](B)

  \node[darkaccent] at (7.65,4.0) {$\alpha$};
  \node at (6.95,1.25) {$a$};
  \node at (7.65,2.9) {$b$};
  \node at (4.0,2.42) {$c$};
\end{tikzpicture}
\end{center}

\key{Почему условие \(a\parallel b\) существенно?}

Если \(a\) и \(b\) пересекаются в одной точке \(Q\), прямая \(c\) может
пересекать обе прямые именно в точке \(Q\) и затем выходить из плоскости.
Тогда известна только одна общая точка \(c\) и \(\alpha\), чего недостаточно,
чтобы утверждать \(c\subset\alpha\).

\medskip

\key{Типичная ошибка:}
увидеть на рисунке, что прямая ``как будто лежит'' в плоскости, и использовать
это без доказательства.

% ----------------------------------------------------------------
\lessonsection{3. Параллельные прямые задают плоскость}

Через две различные параллельные прямые проходит единственная плоскость.

Если
\[
a\parallel b,
\]
то существует единственная плоскость \(\beta\), для которой
\[
a\subset\beta,\qquad b\subset\beta.
\]

Это позволяет вводить \key{вспомогательную плоскость}, даже если она прямо
не названа в условии.

Аналогично единственную плоскость задают две пересекающиеся прямые.

\medskip

\key{Ограничение.}
Две скрещивающиеся прямые одной плоскости не задают.

% ----------------------------------------------------------------
\lessonsection{4. Пересечение двух плоскостей}

Если две различные плоскости имеют общую точку, то их пересечение является
прямой.

Если
\[
\alpha\ne\beta
\]
и у плоскостей найдены две различные общие точки \(A\) и \(B\), то

\[
\boxed{\alpha\cap\beta=AB}.
\]

Отсюда следует особенно важная идея:

\[
\boxed{\text{у двух различных плоскостей не бывает двух разных прямых пересечения}}
\]

Эта мысль позволяет доказывать коллинеарность точек в пространственных
конструкциях.

% ----------------------------------------------------------------
\lessonsection{5. Средняя линия внутри пространственной конструкции}

Пусть \(ABCD\) --- тетраэдр, а точки

\[
M\in DB,\qquad
N\in DC,\qquad
P\in AB,\qquad
Q\in AC
\]

являются серединами соответствующих рёбер.

Рассмотрим отдельные грани.

В треугольнике \(DBC\):
\[
MN\parallel BC,\qquad MN=\frac{BC}{2}.
\]

В треугольнике \(ABC\):
\[
PQ\parallel BC,\qquad PQ=\frac{BC}{2}.
\]

Следовательно,
\[
MN\parallel PQ,\qquad MN=PQ.
\]

Аналогично:

в треугольнике \(DAB\)
\[
MP\parallel AD,\qquad MP=\frac{AD}{2};
\]

в треугольнике \(DCA\)
\[
NQ\parallel AD,\qquad NQ=\frac{AD}{2}.
\]

Поэтому
\[
MP\parallel NQ,\qquad MP=NQ.
\]

Итак, \(MNPQ\) --- параллелограмм, а его периметр равен

\[
\begin{aligned}
P_{MNPQ}
&=
2\left(\frac{BC}{2}+\frac{AD}{2}\right)\\
&=BC+AD.
\end{aligned}
\]

При
\[
AD=12,\qquad BC=14
\]
получаем
\[
\boxed{P_{MNPQ}=26}.
\]

\begin{center}
\begin{tikzpicture}[scale=1.0]
  \tkzDefPoint(0,0){B}
  \tkzDefPoint(6.5,0){C}
  \tkzDefPoint(4.4,4.5){D}
  \tkzDefPoint(2.1,2.4){A}

  \tkzDrawSegments(B,C C,D D,B A,B A,C A,D)

  \coordinate (M) at ($(D)!0.5!(B)$);
  \coordinate (N) at ($(D)!0.5!(C)$);
  \coordinate (P) at ($(A)!0.5!(B)$);
  \coordinate (Q) at ($(A)!0.5!(C)$);

  \draw[darkaccent,line width=1.4pt]
    (M)--(N)--(Q)--(P)--cycle;

  \tkzDrawPoints(B,C,D,A)
  \fill[accent] (M) circle (1.5pt);
  \fill[accent] (N) circle (1.5pt);
  \fill[accent] (P) circle (1.5pt);
  \fill[accent] (Q) circle (1.5pt);

  \tkzLabelPoints[below](B,C)
  \tkzLabelPoints[above](D)
  \tkzLabelPoints[left](A)

  \node[left] at (M) {$M$};
  \node[right] at (N) {$N$};
  \node[left] at (P) {$P$};
  \node[right] at (Q) {$Q$};
\end{tikzpicture}
\end{center}

\key{Полезный приём.}
Если пространственный рисунок мешает увидеть среднюю линию, отдельно
перерисуйте нужную грань как обычный треугольник.

% ----------------------------------------------------------------
\lessonsection{6. Параллельные прямые и подобие треугольников}

Пусть точка \(C\) лежит на отрезке \(AB\). Через точку \(A\) проходит
плоскость \(\alpha\).

Через \(B\) и \(C\) проведены параллельные прямые
\[
BB_1\parallel CC_1,
\]
которые пересекают \(\alpha\) соответственно в точках \(B_1\) и \(C_1\).

Сначала требуется обосновать, что точки
\[
A,\ B_1,\ C_1
\]
лежат на одной прямой.

\subsection*{\color{darkaccent}Почему \(A,B_1,C_1\) коллинеарны}

Прямые \(BB_1\) и \(CC_1\) параллельны, поэтому они задают некоторую
плоскость \(\beta\).

Поскольку
\[
A,B,C
\]
лежат на одной прямой и точки \(B,C\in\beta\), вся прямая \(AB\) лежит
в \(\beta\). Следовательно,
\[
A\in\beta.
\]

Кроме того,
\[
B_1,C_1\in\beta.
\]

По условию
\[
A,B_1,C_1\in\alpha.
\]

Значит, все три точки принадлежат одновременно двум различным плоскостям
\(\alpha\) и \(\beta\). Их пересечение --- одна прямая:

\[
\alpha\cap\beta=AB_1=AC_1.
\]

Следовательно,
\[
\boxed{A,\ B_1,\ C_1\text{ лежат на одной прямой}.}
\]

Теперь существуют два треугольника:
\[
\triangle ACC_1
\qquad\text{и}\qquad
\triangle ABB_1.
\]

Так как
\[
CC_1\parallel BB_1,
\]
эти треугольники подобны:

\[
\triangle ACC_1\sim\triangle ABB_1.
\]

Поэтому
\[
\boxed{
\frac{CC_1}{BB_1}
=
\frac{AC}{AB}
}
\]

и, следовательно,

\[
\boxed{
CC_1=BB_1\cdot\frac{AC}{AB}
}.
\]

\subsection*{\color{darkaccent}Частный случай: \(C\) --- середина \(AB\)}

Если
\[
AC=CB,
\]
то
\[
AC=\frac{AB}{2}.
\]

Следовательно,
\[
CC_1=\frac{BB_1}{2}.
\]

Например, при
\[
BB_1=7
\]
получаем
\[
\boxed{CC_1=3{,}5}.
\]

\subsection*{\color{darkaccent}Отношение \(AC:CB=3:2\)}

Здесь важно сравнивать \(AC\) со всем \(AB\):

\[
AC:CB=3:2
\quad\Longrightarrow\quad
AC=3x,\quad CB=2x,\quad AB=5x.
\]

Следовательно,
\[
\frac{AC}{AB}=\frac35.
\]

Если
\[
BB_1=20,
\]
то
\[
CC_1
=
20\cdot\frac35
=
\boxed{12}.
\]

\key{Типичная ошибка:}
использовать отношение \(CB:AB\), то есть \(\frac25\), хотя стороне
\(CC_1\) в подобии соответствует именно отрезок \(AC\).

% ----------------------------------------------------------------
\lessonsection{7. Как делать пространственный чертёж}

\begin{enumerate}[leftmargin=9mm,itemsep=2mm]
  \item Сначала изобразите основную прямую или плоскость.
  \item Нанесите только те точки, принадлежность которых задана условием.
  \item Если прямая проходит за плоскостью или гранью, скрытый участок
        можно показать штриховой линией.
  \item Не располагайте точку в плоскости только потому, что так рисунок
        выглядит удобнее.
  \item После построения отдельно проверьте каждую связь:
        принадлежность, пересечение, параллельность.
\end{enumerate}

\key{Важно:}
пространственный рисунок является моделью. Он помогает рассуждать, но
доказательство строится на свойствах и теоремах.


% ================================================================
% ИТОГОВЫЙ ЧЕК-ЛИСТ
% ================================================================
\lessonsection{8. Перед тем как считать или доказывать}

Проверьте себя:

\begin{itemize}[leftmargin=8mm,itemsep=2.5mm]
  \item[$\square$] Я различаю обозначения
  \(A\in a\), \(a\subset\alpha\), \(a\parallel b\).
  \item[$\square$] Я могу доказать принадлежность прямой плоскости
  через две её точки.
  \item[$\square$] Я помню, что параллельные прямые задают единственную
  плоскость.
  \item[$\square$] Я умею находить прямую пересечения двух плоскостей
  по двум общим точкам.
  \item[$\square$] Я не считаю факт с рисунка доказанным автоматически.
  \item[$\square$] Я умею выделять нужный треугольник на грани
  пространственной фигуры.
  \item[$\square$] Я узнаю среднюю линию даже внутри тетраэдра.
  \item[$\square$] В подобии я сначала устанавливаю соответствие сторон,
  затем составляю пропорцию.
\end{itemize}

\subsection*{\color{darkaccent}Минимум, который стоит помнить}

\[
\begin{aligned}
&A,B\in c,\quad A,B\in\alpha,\quad A\ne B
&&\Longrightarrow&& c\subset\alpha,\\[2mm]
&a\parallel b
&&\Longrightarrow&&
\text{через }a\text{ и }b\text{ проходит единственная плоскость},\\[2mm]
&\alpha\ne\beta,\quad A,B\in\alpha\cap\beta
&&\Longrightarrow&&
\alpha\cap\beta=AB,\\[2mm]
&M,N\text{ --- середины двух сторон }\triangle ABC
&&\Longrightarrow&&
MN\parallel BC,\quad MN=\frac{BC}{2}.
\end{aligned}
\]

% ================================================================
% ТРЕНИРОВОЧНЫЙ БЛОК
% ================================================================
\clearpage
\thispagestyle{plain}

\begin{center}
{\LARGE\bfseries\color{darkaccent}Тренировочный блок}

\vspace{2mm}

{\large Путь А $\cdot$ 3 задания}
\end{center}

\vspace{7mm}

\textbf{\color{darkaccent}1. Принадлежность прямой плоскости}

В плоскости \(\alpha\) лежат параллельные прямые \(m\) и \(n\).
Прямая \(l\) пересекает \(m\) в точке \(X\), а \(n\) --- в точке \(Y\).

Докажите, что
\[
l\subset\alpha.
\]

\vspace{11mm}

\textbf{\color{darkaccent}2. Средние линии тетраэдра}

В тетраэдре \(ABCD\) точки \(M,N,P,Q\) являются соответственно серединами
рёбер
\[
DB,\quad DC,\quad AB,\quad AC.
\]

Известно:
\[
AD=18,\qquad BC=10.
\]

Найдите периметр четырёхугольника \(MNPQ\).

\vspace{11mm}

\textbf{\color{darkaccent}3. Подобие после построения вспомогательной плоскости}

Точка \(C\) лежит на отрезке \(AB\), причём
\[
AC:CB=2:3.
\]

Через точку \(A\) проведена плоскость \(\alpha\).
Через точки \(B\) и \(C\) проведены параллельные прямые, пересекающие
\(\alpha\) соответственно в точках \(B_1\) и \(C_1\).

Известно:
\[
BB_1=25.
\]

\begin{enumerate}[label=\alph*),leftmargin=8mm,itemsep=2mm]
  \item докажите, что точки \(A,B_1,C_1\) лежат на одной прямой;
  \item найдите \(CC_1\).
\end{enumerate}

% ================================================================
% ОТВЕТЫ
% ================================================================
\clearpage
\thispagestyle{plain}

\begin{center}
{\LARGE\bfseries\color{darkaccent}Ответы}
\end{center}

\vspace{8mm}

\renewcommand{\arraystretch}{1.35}

\begin{table}[ht]
\centering
\caption{Ответы к тренировочному блоку}
\begin{tabularx}{\linewidth}{@{}cX@{}}
\toprule
\textbf{№} & \textbf{Ответ / ключевой результат} \\
\midrule

1 &
Пусть \(X=l\cap m\), \(Y=l\cap n\).
Так как \(m,n\subset\alpha\), имеем \(X,Y\in\alpha\).
Поскольку \(m\parallel n\), точки \(X\ne Y\).
Следовательно, две различные точки прямой \(l\) лежат в \(\alpha\), поэтому
\[
l\subset\alpha.
\]
\\

\addlinespace

2 &
\[
MN=PQ=\frac{BC}{2}=5,
\qquad
MP=NQ=\frac{AD}{2}=9.
\]
Следовательно,
\[
P_{MNPQ}=2(5+9)=\boxed{28}.
\]
\\

\addlinespace

3 &
Прямые \(BB_1\parallel CC_1\) задают вспомогательную плоскость \(\beta\).
Точки \(A,B_1,C_1\) принадлежат одновременно \(\alpha\) и \(\beta\), поэтому
лежат на единственной прямой \(\alpha\cap\beta\).

Далее
\[
\triangle ACC_1\sim\triangle ABB_1.
\]
Так как
\[
AC:CB=2:3,
\]
то
\[
\frac{AC}{AB}=\frac{2}{5}.
\]
Поэтому
\[
CC_1
=
25\cdot\frac25
=
\boxed{10}.
\]
\\

\bottomrule
\end{tabularx}
\end{table}

\vfill

\begin{center}
{\small\color{softgray}
Главная идея: сначала найдите плоскость, в которой задача становится плоской.}
\end{center}

\end{document}